\subsection{Short-form Uncertainty on the LM-Polygraph Benchmark}
\label{sec:polygraph}

The LM-Polygraph benchmark
\citep{fadeeva2023lmpolygraph,vashurin2025lmpolygraph} stresses the
sequence-level use case: each test item is a short response and the
goal is to rank responses by likely hallucination, not to localize
within them. Table~\ref{tab:polygraph} reports per-dataset PRR on
MMLU and GSM8k (QA, binary accuracy) and WMT19 De-En (NMT, COMET
binarised at the training median). The Temp-ViT row aggregates the
Stage~2 per-window logits to one number per response via the same
top-$k$ MIL rule used at training time
(Section~\ref{sec:temporal_head}), so the method is plug-compatible
with the LM-Polygraph sequence-level baselines. This setting answers
(Q\textsubscript{2}): the per-token output Stage~2 was designed to
produce remains competitive with sequence-level scores when one is
forced to commit to a single number per response. Stage 3, which is
trained on long-form passages, shows a small accuracy hit on the
short-form sequences of LM-Polygraph
(Table~\ref{tab:polygraph}); the per-atom loss does not transfer
back to single-claim sequences as cleanly as the Stage 2 MIL
objective.
Table~\ref{tab:polygraph_stats} reports the per-dataset sequence
length distributions.

\begin{table*}[t]
\centering
\caption{\textbf{Long-form hallucination detection --- token level.} We
report AUROC per (LLM, dataset) cell. Models: Mistral-7B-Instruct-v0.2
(Mis-7B-Inst), Llama-3.1-8B-Instruct (LlaMa-8B-Inst),
Llama-3.1-70B-Instruct (LlaMa-70B-Inst). Datasets: RAGTruth
\citep{niu2024ragtruth}, FactScore \citep{min2023}, LongFact
\citep{wei2024}. Per-token gold mask built from atom spans via
longer-span-wins on containment, earlier-span-wins on partial overlap;
uncovered tokens dropped. Best score per column in \textbf{bold},
second-best \underline{underlined}. $^\dagger$ marks cells whose data
comes from a closely related model version
(Mistral-7B-Instruct-v0.1 for RAGTruth, Llama-3-8B-Instruct for
LongFact). Cells with ``---'' indicate experiments not yet run.}
\label{tab:main}
\setlength{\tabcolsep}{2.5pt}
\renewcommand{\arraystretch}{1.05}
\scriptsize
\adjustbox{max width=\textwidth}{%
\begin{tabular}{l ccc ccc ccc}
\toprule
 & \multicolumn{3}{c}{Mis-7B-Inst} & \multicolumn{3}{c}{LlaMa-8B-Inst} & \multicolumn{3}{c}{LlaMa-70B-Inst} \\
\cmidrule(lr){2-4} \cmidrule(lr){5-7} \cmidrule(lr){8-10}
Method & RAGTruth & FactScore & LongFact & RAGTruth & FactScore & LongFact & RAGTruth & FactScore & LongFact \\
\midrule
SemanticEntropy           & --- & --- & --- & --- & --- & --- & --- & --- & --- \\
DegMat (NLI)              & --- & --- & --- & --- & --- & --- & --- & --- & --- \\
EigValLaplacian (NLI)     & --- & --- & --- & --- & --- & --- & --- & --- & --- \\
SAR                       & --- & --- & --- & --- & --- & --- & --- & --- & --- \\
LUQ                       & --- & --- & --- & --- & --- & --- & --- & --- & --- \\
KernelLanguageEntropy     & --- & --- & --- & --- & --- & --- & --- & --- & --- \\
\midrule
Perplexity                & --- & --- & --- & --- & --- & --- & --- & --- & --- \\
MeanTokenEntropy          & --- & --- & --- & --- & --- & --- & --- & --- & --- \\
MCNSE                     & --- & --- & --- & --- & --- & --- & --- & --- & --- \\
TokenSAR                  & --- & --- & --- & --- & --- & --- & --- & --- & --- \\
\midrule
CocoaMSP                  & --- & --- & --- & --- & --- & --- & --- & --- & --- \\
CocoaPPL                  & --- & --- & --- & --- & --- & --- & --- & --- & --- \\
CocoaMTE                  & --- & --- & --- & --- & --- & --- & --- & --- & --- \\
\midrule
Act-ViT (F-v2, N\textsubscript{eff}=6)
                          & \underline{0.570}$^{\dagger}$ & \underline{0.676} & \underline{0.638} & --- & \textbf{0.673} & \underline{0.578}$^{\dagger}$ & --- & --- & --- \\
Temp-ViT (Stage 3, ours)  & \textbf{0.865}$^{\dagger}$ & \textbf{0.696} & \textbf{0.692} & --- & \underline{0.653} & \textbf{0.601}$^{\dagger}$ & --- & --- & --- \\
\bottomrule
\end{tabular}%
}
\end{table*}

\begin{table*}[t]
\centering
\caption{Short-form uncertainty on the {\bf LM-Polygraph benchmark}
\citep{fadeeva2023lmpolygraph,vashurin2025lmpolygraph}. We report
AUROC per dataset over three task families: QA (TriviaQA, CoQA, MMLU,
GSM8k; binary accuracy), NMT (WMT14 Fr--En, WMT19 De--En; COMET
binarised at the training-set median), and summarization (XSUM;
AlignScore binarised at the training-set median). All methods are
evaluated under greedy decoding with the same generations and the same
gold quality measure. Methods are grouped by category:
\emph{Confidence} (sequence-level probability / entropy baselines),
\emph{Mechanistic} (model-internal signal baselines), and \emph{Ours}.
Best score per column is in \textbf{bold}; second-best is
\underline{underlined}. Cells with ``---'' indicate experiments not
yet run.}
\label{tab:polygraph}
\setlength{\tabcolsep}{2pt}
\renewcommand{\arraystretch}{1.05}
\scriptsize
\adjustbox{max width=\textwidth}{%
\begin{tabular}{l *{7}{c} *{7}{c} *{7}{c}}
\toprule
 & \multicolumn{7}{c}{Mis-7B-Inst} & \multicolumn{7}{c}{LlaMa-13B-Inst} & \multicolumn{7}{c}{LlaMa-70B-Inst} \\
\cmidrule(lr){2-8} \cmidrule(lr){9-15} \cmidrule(lr){16-22}
Method
  & TriviaQA & CoQA & MMLU & GSM8k & WMT14 & WMT19 & XSUM
  & TriviaQA & CoQA & MMLU & GSM8k & WMT14 & WMT19 & XSUM
  & TriviaQA & CoQA & MMLU & GSM8k & WMT14 & WMT19 & XSUM \\
\midrule
\multicolumn{22}{l}{\textit{Confidence}} \\
\cmidrule(l){1-22}
MaximumSequenceProbability
  & --- & --- & 0.3563 & 0.7403 & --- & 0.6796 & ---
  & --- & --- & ---    & ---    & --- & ---    & ---
  & --- & --- & ---    & ---    & --- & ---    & --- \\
Perplexity
  & --- & --- & 0.3563 & 0.6981 & --- & 0.7030 & ---
  & --- & --- & ---    & ---    & --- & ---    & ---
  & --- & --- & ---    & ---    & --- & ---    & --- \\
MeanTokenEntropy
  & --- & --- & 0.3554 & 0.7021 & --- & 0.7330 & ---
  & --- & --- & ---    & ---    & --- & ---    & ---
  & --- & --- & ---    & ---    & --- & ---    & --- \\
SelfCertainty
  & --- & --- & 0.3525 & 0.6610 & --- & 0.7766 & ---
  & --- & --- & ---    & ---    & --- & ---    & ---
  & --- & --- & ---    & ---    & --- & ---    & --- \\
TokenSAR
  & --- & --- & 0.5238 & 0.6963 & --- & 0.7013 & ---
  & --- & --- & ---    & ---    & --- & ---    & ---
  & --- & --- & ---    & ---    & --- & ---    & --- \\
MeanPointwiseMutualInformation
  & --- & --- & 0.8072 & 0.4533 & --- & 0.5470 & ---
  & --- & --- & ---    & ---    & --- & ---    & ---
  & --- & --- & ---    & ---    & --- & ---    & --- \\
BoostedProbSequence
  & --- & --- & 0.4072 & 0.7011 & --- & 0.7459 & ---
  & --- & --- & ---    & ---    & --- & ---    & ---
  & --- & --- & ---    & ---    & --- & ---    & --- \\
\midrule
\multicolumn{22}{l}{\textit{Mechanistic}} \\
\cmidrule(l){1-22}
CSL
  & --- & --- & 0.3349 & 0.7015 & --- & 0.6931 & ---
  & --- & --- & ---    & ---    & --- & ---    & ---
  & --- & --- & ---    & ---    & --- & ---    & --- \\
AttentionScore
  & --- & --- & 0.4954 & 0.4920 & --- & 0.5100 & ---
  & --- & --- & ---    & ---    & --- & ---    & ---
  & --- & --- & ---    & ---    & --- & ---    & --- \\
RAUQ
  & --- & --- & 0.4066 & 0.7128 & --- & 0.6561 & ---
  & --- & --- & ---    & ---    & --- & ---    & ---
  & --- & --- & ---    & ---    & --- & ---    & --- \\
FactProbe
  & --- & --- & ---    & ---    & --- & ---    & ---
  & --- & --- & ---    & ---    & --- & ---    & ---
  & --- & --- & ---    & ---    & --- & ---    & --- \\
Feature-Gaps
  & --- & --- & ---    & ---    & --- & ---    & ---
  & --- & --- & ---    & ---    & --- & ---    & ---
  & --- & --- & ---    & ---    & --- & ---    & --- \\
Act-ViT (F-v2, N\textsubscript{eff}=6)
  & --- & --- & \underline{0.9165} & \textbf{0.8643} & --- & 0.7775 & ---
  & --- & --- & ---    & ---    & --- & ---    & ---
  & --- & --- & ---    & ---    & --- & ---    & --- \\
\midrule
\multicolumn{22}{l}{\textit{Ours}} \\
\cmidrule(l){1-22}
Temp-ViT (Encoder Only)
  & --- & --- & 0.9108 & \underline{0.8604} & --- & \textbf{0.7989} & ---
  & --- & --- & ---    & ---    & --- & ---    & ---
  & --- & --- & ---    & ---    & --- & ---    & --- \\
Temp-ViT
  & --- & --- & \textbf{0.9214} & 0.8509 & --- & \underline{0.7835} & ---
  & --- & --- & ---    & ---    & --- & ---    & ---
  & --- & --- & ---    & ---    & --- & ---    & --- \\
\bottomrule
\end{tabular}%
}
\end{table*}

